\documentclass[11pt]{article}

% ------------------------------------------------------------------
% Paquetes
% ------------------------------------------------------------------
\usepackage[utf8]{inputenc}
\usepackage[T1]{fontenc}
% babel español no disponible en este entorno; usamos etiquetas manuales
\usepackage{babel}
\usepackage{amsmath,amssymb,amsthm,mathtools}
\usepackage{bm}
\usepackage[margin=2.5cm]{geometry}
\usepackage{enumitem}
\usepackage{booktabs}
\usepackage{xcolor}
\usepackage{hyperref}
\usepackage{listings}
\usepackage{parskip}
\usepackage{fancyhdr}

\hypersetup{
  colorlinks=true,
  linkcolor=blue!60!black,
  urlcolor=blue!50!black,
  citecolor=blue!60!black
}

% ------------------------------------------------------------------
% Entornos de teoremas
% ------------------------------------------------------------------
\theoremstyle{plain}
\newtheorem{teorema}{Teorema}[section]
\newtheorem{proposicion}[teorema]{Proposición}
\newtheorem{lema}[teorema]{Lema}
\newtheorem{corolario}[teorema]{Corolario}
\theoremstyle{definition}
\newtheorem{definicion}[teorema]{Definición}
\newtheorem{ejemplo}[teorema]{Ejemplo}
\theoremstyle{remark}
\newtheorem{observacion}[teorema]{Observación}

% ------------------------------------------------------------------
% Macros
% ------------------------------------------------------------------
\newcommand{\R}{\mathbb{R}}
\newcommand{\E}{\mathbb{E}}
\newcommand{\Var}{\operatorname{Var}}
\newcommand{\Cov}{\operatorname{Cov}}
\newcommand{\logit}{\operatorname{logit}}
\newcommand{\indep}{\perp\!\!\!\perp}
\newcommand{\bbeta}{\bm{\beta}}
\newcommand{\bx}{\bm{x}}
\newcommand{\bX}{\bm{X}}
\newcommand{\by}{\bm{y}}
\newcommand{\bu}{\bm{u}}
\newcommand{\bW}{\bm{W}}
\newcommand{\bI}{\bm{I}}
\newcommand{\bH}{\bm{H}}
\newcommand{\btheta}{\bm{\theta}}
\newcommand{\bpi}{\bm{\pi}}
\newcommand{\hbeta}{\hat{\bm{\beta}}}
\newcommand{\convd}{\xrightarrow{d}}
\newcommand{\convp}{\xrightarrow{p}}
\newcommand{\iid}{\overset{\text{iid}}{\sim}}
\newcommand{\Bern}{\operatorname{Bernoulli}}
\newcommand{\Nor}{\mathcal{N}}

% ------------------------------------------------------------------
% Código Python
% ------------------------------------------------------------------
\definecolor{codebg}{rgb}{0.96,0.96,0.96}
\definecolor{codekw}{rgb}{0.55,0.0,0.55}
\definecolor{codestr}{rgb}{0.0,0.45,0.0}
\definecolor{codecom}{rgb}{0.4,0.4,0.4}
\lstdefinestyle{py}{
  language=Python,
  backgroundcolor=\color{codebg},
  basicstyle=\ttfamily\footnotesize,
  keywordstyle=\color{codekw}\bfseries,
  stringstyle=\color{codestr},
  commentstyle=\color{codecom}\itshape,
  showstringspaces=false,
  breaklines=true,
  frame=single,
  rulecolor=\color{gray!40},
  numbers=left,
  numberstyle=\tiny\color{gray},
  xleftmargin=2em,
  framexleftmargin=1.5em
}

\pagestyle{fancy}
\fancyhf{}
\rhead{Regresión logística: teoría y diagnóstico}
\lhead{}
\cfoot{\thepage}

% ==================================================================
\title{\textbf{Regresión Logística:\\
Estimación por Máxima Verosimilitud, Inferencia Asintótica\\
y Diagnóstico de Linealidad en el Logit}}
\author{Notas de estudio --- nivel ISLR / teoría de verosimilitud}
\date{\today}

\begin{document}
\maketitle
\tableofcontents
\newpage

% ==================================================================
\section{Introducción y objetivo}
% ==================================================================

Estas notas desarrollan, desde primeros principios, la maquinaria inferencial
de la regresión logística y los tests de diagnóstico que verifican el supuesto
central del modelo: la \emph{linealidad en el logit}. El estándar de rigor es
el de un curso de teoría de la verosimilitud: para cada estadístico se deriva
su forma, se establece su distribución bajo $H_0$ y se justifica por qué se usa
ese y no otro.

El hilo conductor es el siguiente flujo de trabajo, que ejecutaremos sobre un
problema real de \emph{scoring} de crédito (dataset \emph{Home Credit Default
Risk}):

\begin{enumerate}[label=\textbf{\arabic*.}]
  \item Ajuste del modelo logístico por máxima verosimilitud (MLE), vía
        Newton--Raphson / IRLS.
  \item Diagnóstico visual de linealidad en el logit (\emph{empirical logit
        plot}).
  \item Test formal de Box--Tidwell (basado en el test de Wald sobre un término
        construido).
  \item Remedios ante violaciones (transformaciones, términos polinomiales).
  \item Bondad de ajuste global (test de Hosmer--Lemeshow).
\end{enumerate}

La \textbf{trinidad} de tests asintóticos ---Wald, Score (Rao) y Razón de
Verosimilitudes (LRT)--- aparece transversalmente, de modo que la Sección
\ref{sec:trinidad} la desarrolla en general y las secciones siguientes la
instancian.

% ==================================================================
\section{El modelo logístico}\label{sec:modelo}
% ==================================================================

\subsection{Planteamiento y motivación del enlace}

Sea $Y \in \{0,1\}$ una respuesta binaria y $\bx = (1, x_1, \dots, x_p)^\top
\in \R^{p+1}$ un vector de predictores (con intercepto). Modelamos
\[
  Y_i \mid \bx_i \iid \Bern(\pi_i), \qquad
  \pi_i = \Pr(Y_i = 1 \mid \bx_i).
\]

Un modelo lineal ingenuo $\pi_i = \bx_i^\top\bbeta$ falla por dos razones,
ambas consecuencia de que $Y$ es Bernoulli:

\begin{itemize}
  \item \textbf{Rango:} $\bx_i^\top\bbeta$ recorre todo $\R$, pero
        $\pi_i \in (0,1)$. Un modelo lineal produce probabilidades fuera de
        $[0,1]$.
  \item \textbf{Heterocedasticidad:} $\Var(Y_i) = \pi_i(1-\pi_i)$ depende de
        $\pi_i$ y por tanto de $\bx_i$; no es constante. Esto viola el supuesto
        de homocedasticidad del modelo lineal (Gauss--Markov), inutilizando OLS.
\end{itemize}

La solución es introducir una \textbf{función de enlace} $g$ que mapee
$(0,1) \to \R$. El enlace \emph{canónico} para la familia Bernoulli es el
\textbf{logit}:
\[
  \boxed{\;\logit(\pi_i)
  = \log\!\frac{\pi_i}{1-\pi_i}
  = \bx_i^\top \bbeta \;}
\]
La cantidad $\pi/(1-\pi)$ son los \emph{odds} (momios), y $\log\{\pi/(1-\pi)\}$
el \emph{log-odds}. Invirtiendo:
\[
  \pi_i = \frac{e^{\bx_i^\top\bbeta}}{1 + e^{\bx_i^\top\bbeta}}
  = \frac{1}{1 + e^{-\bx_i^\top\bbeta}}
  = \sigma(\bx_i^\top\bbeta),
\]
donde $\sigma(\cdot)$ es la función logística (sigmoide).

\begin{observacion}[¿Por qué el logit y no otro enlace?]
El logit es el \textbf{enlace canónico} de la familia exponencial Bernoulli.
Escribiendo la pmf de Bernoulli en forma canónica,
\[
  f(y;\pi) = \pi^y(1-\pi)^{1-y}
  = \exp\!\Big\{ y \underbrace{\log\tfrac{\pi}{1-\pi}}_{\theta}
  + \log(1-\pi) \Big\},
\]
el parámetro natural es exactamente $\theta = \logit(\pi)$. Modelar
$\theta = \bx^\top\bbeta$ da simplificaciones notables: la información observada
coincide con la esperada, y Newton--Raphson coincide con Fisher scoring
(Sección \ref{sec:irls}). Otros enlaces son válidos (probit: $g=\Phi^{-1}$;
complementario log-log) pero no canónicos.
\end{observacion}

\subsection{Interpretación de los coeficientes}

De $\logit(\pi) = \bx^\top\bbeta$, un incremento unitario en $x_j$ (con lo demás
fijo) cambia el log-odds en $\beta_j$, es decir multiplica los \emph{odds} por
$e^{\beta_j}$. Por eso $e^{\beta_j}$ se llama \emph{odds ratio}. Este carácter
\textbf{multiplicativo} (frente al aditivo de la regresión lineal) es esencial
al interpretar variables transformadas: si el modelo usa $\log(x_j)$, entonces
$\beta_j$ mide el efecto de multiplicar $x_j$ por $e$.

% ==================================================================
\section{Estimación por máxima verosimilitud}\label{sec:mle}
% ==================================================================

A diferencia de la regresión lineal (donde OLS da forma cerrada), en logística
no hay solución cerrada y se recurre al MLE resuelto numéricamente. Seguimos la
teoría de verosimilitud estándar (cf. Rodríguez, Princeton, \emph{GLM Notes},
Ap.~A).

\subsection{Log-verosimilitud}

Como los $Y_i$ son Bernoulli independientes, la verosimilitud es
\[
  L(\bbeta) = \prod_{i=1}^n \pi_i^{y_i}(1-\pi_i)^{1-y_i},
\]
y la \textbf{log-verosimilitud}:
\begin{align}
  \ell(\bbeta)
  &= \sum_{i=1}^n \big[ y_i \log \pi_i + (1-y_i)\log(1-\pi_i) \big] \notag\\
  &= \sum_{i=1}^n \big[ y_i \log\tfrac{\pi_i}{1-\pi_i} + \log(1-\pi_i)\big]
     \notag\\
  &= \sum_{i=1}^n \big[ y_i\, \bx_i^\top\bbeta
     - \log\!\big(1 + e^{\bx_i^\top\bbeta}\big) \big],
  \label{eq:loglik}
\end{align}
donde en la última línea usamos $\logit(\pi_i)=\bx_i^\top\bbeta$ y
$\log(1-\pi_i) = -\log(1+e^{\bx_i^\top\bbeta})$.

\subsection{Vector score}

El \textbf{score} es el gradiente de $\ell$. Usando
$\dfrac{\partial}{\partial\bbeta}\log(1+e^{\bx_i^\top\bbeta})
= \dfrac{e^{\bx_i^\top\bbeta}}{1+e^{\bx_i^\top\bbeta}}\bx_i = \pi_i \bx_i$,
\begin{equation}
  \bu(\bbeta)
  = \frac{\partial \ell}{\partial \bbeta}
  = \sum_{i=1}^n (y_i - \pi_i)\,\bx_i
  = \bX^\top(\by - \bpi),
  \label{eq:score}
\end{equation}
donde $\bX$ es la matriz de diseño $n\times(p+1)$, $\by=(y_i)$ y
$\bpi=(\pi_i)$. Las \emph{ecuaciones de verosimilitud} $\bu(\hbeta)=\bm 0$
igualan los momentos observados y ajustados: $\bX^\top\by = \bX^\top\hat\bpi$.

\begin{observacion}[Media cero del score]
Como $\E[Y_i]=\pi_i$, se tiene $\E[\bu(\bbeta)] = \bX^\top(\E[\by]-\bpi)=\bm 0$
en el valor verdadero. Esta es la primera identidad de Bartlett, base de toda
la teoría asintótica.
\end{observacion}

\subsection{Hessiano e información de Fisher}

Derivando \eqref{eq:score} otra vez, y usando
$\partial \pi_i/\partial\bbeta = \pi_i(1-\pi_i)\bx_i$:
\begin{equation}
  \bH(\bbeta)
  = \frac{\partial^2 \ell}{\partial\bbeta\,\partial\bbeta^\top}
  = -\sum_{i=1}^n \pi_i(1-\pi_i)\,\bx_i\bx_i^\top
  = -\bX^\top \bW \bX,
  \label{eq:hessian}
\end{equation}
donde $\bW = \operatorname{diag}\big(\pi_i(1-\pi_i)\big)$. Como $\bH$ no depende
de $\by$ (solo de $\bpi$, que depende de $\bbeta$ y $\bX$), la
\textbf{información observada} $-\bH$ coincide con la \textbf{información
esperada} (de Fisher):
\begin{equation}
  \bI(\bbeta) = \E[-\bH(\bbeta)] = \bX^\top\bW\bX.
  \label{eq:fisher}
\end{equation}
Esta coincidencia es precisamente la propiedad del enlace canónico. Además
$\bI(\bbeta)$ es definida positiva (si $\bX$ tiene rango completo y
$0<\pi_i<1$), luego $\ell$ es cóncava y el MLE, cuando existe, es único.

\subsection{Newton--Raphson e IRLS}\label{sec:irls}

Como $\bu(\hbeta)=\bm 0$ no tiene forma cerrada, expandimos el score en serie
de Taylor de primer orden alrededor de un valor $\bbeta^{(t)}$:
\[
  \bm 0 = \bu(\hbeta) \approx \bu(\bbeta^{(t)})
  + \bH(\bbeta^{(t)})\,(\hbeta - \bbeta^{(t)}),
\]
y despejando obtenemos la iteración de \textbf{Newton--Raphson}:
\begin{equation}
  \bbeta^{(t+1)}
  = \bbeta^{(t)} - \bH^{-1}(\bbeta^{(t)})\,\bu(\bbeta^{(t)})
  = \bbeta^{(t)} + (\bX^\top\bW\bX)^{-1}\bX^\top(\by-\bpi).
  \label{eq:newton}
\end{equation}
Al usar el enlace canónico, $\bH=-\bI$, por lo que Newton--Raphson coincide con
\textbf{Fisher scoring}. La ecuación \eqref{eq:newton} puede reescribirse como
un problema de mínimos cuadrados ponderados: definiendo la
\emph{variable de respuesta ajustada}
\[
  z_i^{(t)} = \bx_i^\top\bbeta^{(t)}
  + \frac{y_i - \pi_i^{(t)}}{\pi_i^{(t)}(1-\pi_i^{(t)})},
\]
se obtiene
\[
  \bbeta^{(t+1)} = (\bX^\top\bW\bX)^{-1}\bX^\top\bW\bm z^{(t)},
\]
que es exactamente el estimador de mínimos cuadrados ponderados (WLS) de
$\bm z^{(t)}$ sobre $\bX$ con pesos $\bW$. De ahí el nombre \textbf{IRLS}
(\emph{Iteratively Reweighted Least Squares}): en cada paso resolvemos una
regresión lineal ponderada, actualizamos $\bpi$ y $\bW$, y repetimos hasta
convergencia. Este es el algoritmo que implementan \texttt{glm} en R y
\texttt{statsmodels}/\texttt{sklearn} en Python.

\subsection{Distribución asintótica del MLE}

\begin{teorema}[Normalidad asintótica del MLE]\label{thm:asymp}
Bajo condiciones de regularidad (parámetro verdadero interior al espacio,
$\ell$ tres veces diferenciable con terceras derivadas acotadas), el MLE
satisface
\[
  \hbeta \;\overset{a}{\sim}\; \Nor_{p+1}\!\big(\bbeta,\; \bI^{-1}(\bbeta)\big),
  \qquad\text{equivalentemente}\qquad
  \sqrt{n}(\hbeta-\bbeta)\convd \Nor\!\big(\bm 0,\, [n^{-1}\bI(\bbeta)]^{-1}\big).
\]
\end{teorema}

\emph{Idea de la demostración.} Se expande $\bu(\hbeta)=\bm 0$ alrededor de
$\bbeta$: $\bm0 = \bu(\bbeta) + \bH(\bbeta)(\hbeta-\bbeta) + o_p$. Entonces
$\sqrt n(\hbeta-\bbeta) = [-n^{-1}\bH]^{-1}\, n^{-1/2}\bu(\bbeta)+o_p$. Por la
ley de los grandes números $-n^{-1}\bH \convp n^{-1}\bI$, y por el TCL aplicado
a $\bu(\bbeta)=\sum_i(y_i-\pi_i)\bx_i$ (suma de términos independientes de media
cero y varianza $\bI$), $n^{-1/2}\bu(\bbeta)\convd\Nor(\bm0, n^{-1}\bI)$.
Slutsky cierra el argumento. $\qquad\blacksquare$

En la práctica estimamos la covarianza por $\widehat{\Var}(\hbeta) =
\bI^{-1}(\hbeta) = (\bX^\top\hat\bW\bX)^{-1}$, y el \textbf{error estándar} de
$\hat\beta_j$ es
\[
  \operatorname{SE}(\hat\beta_j)
  = \sqrt{\big[(\bX^\top\hat\bW\bX)^{-1}\big]_{jj}}.
\]

\begin{observacion}[Qué es SE]
El \emph{error estándar} $\operatorname{SE}(\hat\beta_j)$ es la desviación
estándar \emph{de la distribución muestral del estimador} $\hat\beta_j$ (no de
los datos). Cuantifica cuánto variaría $\hat\beta_j$ entre muestras repetidas.
Sale de la raíz del elemento diagonal $j$ de la matriz de covarianza asintótica
$\bI^{-1}(\hbeta)$.
\end{observacion}

% ==================================================================
\section{La trinidad de tests asintóticos}\label{sec:trinidad}
% ==================================================================

Sea $H_0:\btheta=\btheta_0$ (o, más generalmente, restricciones sobre un
subvector). Existen tres estadísticos clásicos, todos con distribución
asintótica $\chi^2$ bajo $H_0$. Difieren en \emph{dónde} evalúan la
verosimilitud.

\begin{center}
\begin{tabular}{lll}
\toprule
Test & Evalúa en & Requiere ajustar \\
\midrule
Wald  & $\hbeta$ (irrestricto)      & solo modelo grande \\
Score (Rao) & $\btheta_0$ (restringido) & solo modelo chico \\
LRT   & ambos                       & ambos modelos \\
\bottomrule
\end{tabular}
\end{center}

\subsection{Test de Wald}

\begin{teorema}[Wald]
Bajo $H_0:\btheta=\btheta_0$ y las condiciones del Teorema \ref{thm:asymp}, el
estadístico
\[
  W = (\hbeta-\btheta_0)^\top \bI(\hbeta)\,(\hbeta-\btheta_0)
  \convd \chi^2_r,
\]
donde $r$ es el número de restricciones.
\end{teorema}

\emph{Derivación.} Del Teorema \ref{thm:asymp}, $\hbeta \overset{a}{\sim}
\Nor(\btheta_0, \bI^{-1})$ bajo $H_0$. Estandarizando, la forma cuadrática de un
vector normal $\bm z\sim\Nor(\bm0,\Sigma)$ cumple $\bm z^\top\Sigma^{-1}\bm z
\sim\chi^2_r$. Tomando $\bm z=\hbeta-\btheta_0$ y $\Sigma=\bI^{-1}$ se obtiene
el resultado. $\quad\blacksquare$

Para una sola componente ($r=1$), se toma la raíz con signo:
\[
  z_j = \frac{\hat\beta_j - \beta_{j,0}}{\operatorname{SE}(\hat\beta_j)}
  \convd \Nor(0,1),
  \qquad z_j^2 \convd \chi^2_1.
\]
Este es exactamente el \emph{z-value} que reporta \texttt{statsmodels} en la
tabla de coeficientes. Nótese el paralelo con el test $t$ de regresión lineal:
allí el estadístico sigue una $t$ exacta (por normalidad de los errores); aquí
solo tenemos normalidad \emph{asintótica}, así que la referencia es $\Nor(0,1)$
y no $t$.

\subsection{Test de Score (Rao)}

\begin{teorema}[Score]
Bajo $H_0:\btheta=\btheta_0$, el score satisface
$\bu(\btheta_0)\overset{a}{\sim}\Nor(\bm0,\bI(\btheta_0))$, y el estadístico
\[
  Q = \bu(\btheta_0)^\top \bI^{-1}(\btheta_0)\, \bu(\btheta_0) \convd \chi^2_r.
\]
\end{teorema}

\emph{Derivación.} Bajo $H_0$, $\bu(\btheta_0)=\sum_i(y_i-\pi_{i,0})\bx_i$ es
suma de términos independientes de media cero (identidad de Bartlett) y
covarianza $\bI(\btheta_0)$. Por el TCL,
$\bu(\btheta_0)\convd\Nor(\bm0,\bI(\btheta_0))$. La forma cuadrática
correspondiente es $\chi^2_r$. $\quad\blacksquare$

La gran ventaja: \emph{no requiere calcular el MLE del modelo grande}, solo el
del modelo restringido (a menudo trivial). Por eso es el motor de muchos tests
de diagnóstico (incluido el enfoque de "variable construida" del Box--Tidwell).

\subsection{Test de Razón de Verosimilitudes (LRT)}

\begin{teorema}[Wilks]
Sean $\omega_1\subset\omega_2$ modelos anidados. Bajo $H_0$ (el modelo pequeño
es correcto),
\[
  G^2 = -2\log\lambda
  = 2\big[\ell(\hbeta_{\omega_2}) - \ell(\hbeta_{\omega_1})\big]
  \convd \chi^2_\nu,
  \qquad \nu = \dim(\omega_2)-\dim(\omega_1).
\]
\end{teorema}

\emph{Derivación (esbozo).} Expandiendo $\ell(\hbeta_{\omega_1})$ alrededor de
$\hbeta_{\omega_2}$ en Taylor de segundo orden y usando que el score se anula en
el MLE, se obtiene
$G^2 \approx (\hbeta_{\omega_2}-\hbeta_{\omega_1})^\top
\bI (\hbeta_{\omega_2}-\hbeta_{\omega_1})$,
que es asintóticamente equivalente al estadístico de Wald y por tanto
$\chi^2_\nu$. $\quad\blacksquare$

\begin{observacion}[Equivalencia asintótica y elección]
Los tres estadísticos convergen a la misma $\chi^2_r$ bajo $H_0$ y son
asintóticamente equivalentes. En muestras finitas difieren; simulaciones
sugieren que el LRT suele ser preferible. La elección práctica se rige por el
costo computacional (¿ya tengo ajustado el modelo grande, el chico, o ambos?) y
por robustez. En muestras muy grandes ---como el dataset de crédito con
$n\approx 3\times10^5$--- los tres coinciden casi exactamente.
\end{observacion}

% ==================================================================
\section{Diagnóstico visual: el \emph{empirical logit plot}}\label{sec:visual}
% ==================================================================

El supuesto central de la regresión logística no es que $X$ se relacione
linealmente con $Y$ (para eso está la curva sigmoide), sino que el
\textbf{logit} sea lineal en cada predictor continuo:
\[
  \logit(\pi(x)) = \beta_0 + \beta_1 x .
\]

\textbf{Idea del diagnóstico.} Si agrupamos las observaciones en bins según $x$
y en cada bin estimamos $\hat p_b = \bar y_b$ (proporción de éxitos), entonces
bajo el modelo el logit empírico $\log\{\hat p_b/(1-\hat p_b)\}$ debe ser
aproximadamente lineal en el centro $\bar x_b$ del bin.

\paragraph{Por qué bins de igual frecuencia (\texttt{qcut}).} La varianza del
logit empírico en un bin depende del número de observaciones:
\[
  \Var\big(\logit(\hat p_b)\big) \approx
  \frac{1}{n_b\,\hat p_b(1-\hat p_b)}
\]
(por el método delta aplicado a $\logit$). Bins con pocas observaciones dan
puntos muy ruidosos. Usar cuantiles (\texttt{pd.qcut}) garantiza $n_b$
aproximadamente constante y estabiliza la varianza; usar bins de igual ancho
(\texttt{pd.cut}) dejaría bins casi vacíos en las colas.

\begin{lstlisting}[style=py, caption={Empirical logit plot con manejo de NaN e infinitos.}]
import numpy as np
import pandas as pd
import matplotlib.pyplot as plt

def empirical_logit_plot(X_col, y, n_bins=10):
    df = pd.DataFrame({"X": X_col, "y": y})
    # Robustez: eliminar NaN e infinitos (p.ej. log de valores <= 0)
    df = df.replace([np.inf, -np.inf], np.nan).dropna()

    # Bins de igual frecuencia; duplicates='drop' evita errores por empates
    df["bin"] = pd.qcut(df["X"], q=n_bins, duplicates="drop")
    grouped = df.groupby("bin", observed=True).agg(
        p=("y", "mean"),
        X_mid=("X", "mean"),
    )
    # Solo bins con p en (0,1): el logit no existe en 0 ni 1
    grouped = grouped[(grouped["p"] > 0) & (grouped["p"] < 1)]
    grouped["logit"] = np.log(grouped["p"] / (1 - grouped["p"]))

    plt.figure()
    plt.scatter(grouped["X_mid"], grouped["logit"])
    coef = np.polyfit(grouped["X_mid"], grouped["logit"], 1)
    xs = np.unique(grouped["X_mid"])
    plt.plot(xs, np.poly1d(coef)(xs), "r--")
    plt.xlabel(X_col.name); plt.ylabel("logit empirico")
    plt.title(f"Logit empirico vs {X_col.name}")
    plt.show()
\end{lstlisting}

\paragraph{Limitación.} Con $n$ grande y solo 10 bins, cada punto resume
decenas de miles de observaciones y el diagrama puede verse ruidoso aun cuando
la relación sea lineal. El visual es \emph{preliminar}; la decisión formal la da
el test de Box--Tidwell (Sección \ref{sec:boxtidwell}).

% ==================================================================
\section{Test de Box--Tidwell}\label{sec:boxtidwell}
% ==================================================================

\subsection{Origen: transformación de Box y Tidwell (1962)}

Box y Tidwell (1962, \emph{Technometrics} \textbf{4}, 531--550) propusieron
estimar la potencia $\lambda$ de una transformación $x\mapsto x^\lambda$ que
linealiza la relación. En el contexto logístico, supongamos que la forma
verdadera es
\[
  \logit(\pi) = \beta_0 + \beta_1\, x^\lambda,
\]
y queremos contrastar $H_0:\lambda=1$ (linealidad). El problema es que $\lambda$
entra de forma \emph{no lineal} en los parámetros, lo que complicaría la
estimación. La solución de Box--Tidwell es \textbf{linealizar por Taylor}.

\subsection{Linealización por serie de Taylor}

Expandimos $x^\lambda$ como función de $\lambda$ alrededor de $\lambda=1$.
Recordando que $\dfrac{d}{d\lambda}x^\lambda = x^\lambda\log x$, evaluado en
$\lambda=1$ da $x\log x$:
\[
  x^\lambda \approx x^1 + (\lambda-1)\,\big[x\log x\big]
  = x + (\lambda-1)\,x\log x.
\]
Sustituyendo en el modelo:
\[
  \logit(\pi) \approx \beta_0 + \beta_1 x
  + \underbrace{\beta_1(\lambda-1)}_{\displaystyle \gamma}\, (x\log x).
\]
Definiendo $\gamma = \beta_1(\lambda-1)$, el modelo se vuelve \textbf{lineal en
los parámetros} $(\beta_0,\beta_1,\gamma)$, con un \emph{término construido}
adicional $x\log x$:
\begin{equation}
  \boxed{\;\logit(\pi) = \beta_0 + \beta_1 x + \gamma\,(x\log x)\;}
  \label{eq:bt}
\end{equation}

\subsection{El contraste y su estadístico}

Bajo esta parametrización, la hipótesis de linealidad $\lambda=1$ equivale a
\[
  H_0:\ \gamma = 0 \qquad\text{vs}\qquad H_a:\ \gamma\neq 0.
\]
Como \eqref{eq:bt} es un modelo logístico ordinario (con $x\log x$ como un
predictor más), el contraste sobre $\gamma$ es un \textbf{test de Wald} estándar
sobre un coeficiente individual:
\[
  z = \frac{\hat\gamma}{\operatorname{SE}(\hat\gamma)} \convd \Nor(0,1)
  \quad\text{bajo }H_0,
  \qquad z^2\convd\chi^2_1.
\]
Rechazar $H_0$ ($p<0.05$) indica que el término $x\log x$ mejora
significativamente el ajuste, es decir, \textbf{hay violación de la linealidad
en el logit}. Equivalentemente puede usarse el LRT comparando \eqref{eq:bt} con
el modelo sin $x\log x$; ambos son asintóticamente $\chi^2_1$.

\begin{observacion}[¿Por qué este estadístico y no un test directo sobre
$\lambda$?]
Un test de Wald/LRT directo sobre $\lambda$ exigiría estimar el modelo no lineal
$\beta_0+\beta_1 x^\lambda$, con optimización más costosa e inestable. El truco
de Box--Tidwell convierte el problema en \emph{añadir una covariable
construida} y testear su coeficiente con la maquinaria ya disponible. El costo:
es una prueba \emph{local} (válida cerca de $\lambda=1$) y de primer orden.
\end{observacion}

\subsection{Un test o muchos: Wald escalar, Wald multivariado y LRT}
\label{sec:escalar-vs-conjunto}

Al ajustar el \textbf{único} modelo aumentado con todos los términos
construidos,
\[
  \logit(\pi) = \beta_0 + \sum_{j=1}^m \beta_j x_j
  + \sum_{j=1}^m \gamma_j\,(x_j\log x_j),
\]
disponemos de un vector $\bm\gamma=(\gamma_1,\dots,\gamma_m)^\top$. Sobre este
\emph{mismo} modelo pueden plantearse dos hipótesis distintas, y conviene no
confundirlas.

\paragraph{(A) Hipótesis univariada (una por variable).}
\[
  H_0^{(j)}:\ \gamma_j = 0 \qquad (r=1).
\]
Es el \textbf{Wald escalar} sobre un coeficiente individual, ya derivado en la
Sección \ref{sec:trinidad}:
\[
  z_j = \frac{\hat\gamma_j}{\operatorname{SE}(\hat\gamma_j)}\convd\Nor(0,1),
  \qquad z_j^2\convd\chi^2_1,
\]
donde $\operatorname{SE}(\hat\gamma_j)=\sqrt{[\bI^{-1}(\hbeta)]_{jj}}$. Testea la
linealidad de \emph{una} variable, controlando por todo lo demás (incluidos los
otros $\gamma_k$). Es el que reporta \texttt{statsmodels} en
\texttt{.pvalues}, y el que guía la decisión de \emph{cuál} variable
transformar.

\paragraph{(B) Hipótesis conjunta (todas a la vez).}
\[
  H_0:\ \gamma_1=\gamma_2=\cdots=\gamma_m=0 \qquad (r=m).
\]
La hipótesis es que \emph{ninguna} variable viola linealidad simultáneamente.
Aquí sí aparece un \textbf{Wald multivariado}. Escribiendo la restricción en
forma matricial $H_0:\bm C\hbeta = \bm 0$, donde $\bm C$ es la matriz
$m\times(2m+1)$ que selecciona las componentes $\bm\gamma$ del vector completo,
el estadístico es la \textbf{forma cuadrática}
\begin{equation}
  W = \hat{\bm\gamma}^\top
  \big[\widehat{\Var}(\hat{\bm\gamma})\big]^{-1}\hat{\bm\gamma}
  \convd \chi^2_m,
  \label{eq:wald-multi}
\end{equation}
con $\widehat{\Var}(\hat{\bm\gamma})$ el bloque $m\times m$ de la covarianza
$\bI^{-1}(\hbeta)=(\bX^\top\bW\bX)^{-1}$ correspondiente a $\bm\gamma$.

\emph{Derivación.} Por el Teorema \ref{thm:asymp}, $\hbeta\overset{a}{\sim}
\Nor(\bbeta,\bI^{-1})$. Una transformación lineal de un vector normal es normal:
$\bm C\hbeta \overset{a}{\sim}\Nor(\bm C\bbeta,\ \bm C\bI^{-1}\bm C^\top)$. Bajo
$H_0$, $\bm C\bbeta=\bm 0$, y como $\bm C$ selecciona el subvector $\bm\gamma$,
$\bm C\bI^{-1}\bm C^\top = \widehat{\Var}(\hat{\bm\gamma})$ es justo el bloque
diagonal. La forma cuadrática de un normal $\bm z\sim\Nor(\bm0,\Sigma)$ cumple
$\bm z^\top\Sigma^{-1}\bm z\sim\chi^2_m$; con $\bm z=\hat{\bm\gamma}$ se obtiene
\eqref{eq:wald-multi}. $\quad\blacksquare$

Este es el análogo logístico exacto del estadístico $F$ que en regresión lineal
contrasta un \emph{bloque} de coeficientes simultáneamente (allí $F$ exacta por
normalidad de los errores; aquí $\chi^2_m$ solo asintótica).

\paragraph{El LRT conjunto: alternativa equivalente.}
La misma hipótesis conjunta $H_0:\bm\gamma=\bm0$ puede contrastarse por razón de
verosimilitudes, comparando el modelo aumentado (grande) con el modelo sin los
$m$ términos construidos (chico):
\[
  G^2 = 2\big[\ell(\text{aumentado}) - \ell(\text{reducido})\big]
  \convd\chi^2_m.
\]
Por la equivalencia asintótica de la trinidad (Sección \ref{sec:trinidad}),
$W$ de \eqref{eq:wald-multi} y $G^2$ convergen a la misma $\chi^2_m$ bajo $H_0$.
En la implementación de la Sección \ref{sec:boxtidwell} usamos el \textbf{LRT}
para la parte conjunta ---no el Wald multivariado--- simplemente porque ya
teníamos ambos modelos ajustados y resulta más económico; ambas respuestas son
válidas e intercambiables.

\begin{observacion}[Escalar y conjunto no son redundantes]
Puede ocurrir que \emph{ningún} Wald escalar rechace ($p_j>0.05$ para todo $j$)
pero el test conjunto \emph{sí} rechace: violaciones pequeñas y correlacionadas
suman evidencia. Y a la inversa, un escalar aislado puede rechazar sin que el
conjunto lo haga. Es el mismo fenómeno que en regresión lineal, donde el $F$
global puede ser significativo con todos los $t$ individuales no significativos
(multicolinealidad). La forma cuadrática \eqref{eq:wald-multi} captura la
covarianza \emph{entre} los $\hat\gamma_j$ vía el bloque
$\widehat{\Var}(\hat{\bm\gamma})$, información que los tests escalares uno a uno
ignoran.
\end{observacion}

\begin{center}
\begin{tabular}{lllc}
\toprule
Pregunta & Hipótesis & Test & Distr.\ bajo $H_0$ \\
\midrule
¿\emph{Esta} variable viola linealidad? & $\gamma_j=0$
  & Wald escalar & $\chi^2_1$ \\
¿\emph{Alguna} variable viola linealidad? & $\bm\gamma=\bm0$
  & Wald multivariado \emph{o} LRT & $\chi^2_m$ \\
\bottomrule
\end{tabular}
\end{center}

Para el flujo de decisión (transformar variable por variable) se necesita el
\textbf{Wald escalar}, porque interesa saber \emph{cuál} corregir. El test
conjunto sirve de chequeo global: si el $\chi^2_m$ no rechaza, puede detenerse el
análisis sin revisar variable por variable.

\paragraph{Wald multivariado explícito en Python.}
\begin{lstlisting}[style=py]
# Sobre el modelo aumentado m_full (con los terminos X*ln(X)):
terminos = [c for c in m_full.params.index if c.endswith("_ln")]

# H0 conjunta: todos los gamma = 0  -> Wald multivariado (chi^2 con m gl)
wald_conjunto = m_full.wald_test(
    ", ".join(f"{t} = 0" for t in terminos),
    scalar=True,
)
print(wald_conjunto)   # estadistico chi^2, gl = m, y su p-valor
\end{lstlisting}
Internamente \texttt{wald\_test} construye $\bm C$, extrae el bloque de
covarianza y evalúa la forma cuadrática \eqref{eq:wald-multi}.

\subsection{Restricciones de aplicabilidad}

El término $x\log x$ requiere $x>0$. Por tanto:
\begin{itemize}
  \item \textbf{Continuas positivas:} aplica directamente.
  \item \textbf{Con ceros/negativos:} desplazar ($x\mapsto x-\min+c$) o usar la
        versión con enlace/otra transformación. Si el cero es
        \emph{estructural} (p.ej.\ \texttt{amt\_goods\_price}=0 para préstamos en
        efectivo), la práctica correcta es separar en un flag binario más la
        variable transformada donde $x>0$.
  \item \textbf{Binarias / dummies:} \textbf{no aplica}. Con solo dos valores,
        dos puntos siempre determinan una recta; no hay curva que linealizar, y
        $\log 0$ no existe. El supuesto de linealidad en el logit es un supuesto
        \emph{sobre variables continuas}.
\end{itemize}

\subsection{Por qué incluir el resto de covariables (control)}

El test debe correrse en el \textbf{modelo completo}: se incluyen todos los
predictores y se añade $x_j\log x_j$ para cada continua a testear. Omitir
covariables relevantes induce el \textbf{sesgo por variable omitida} (misma
lógica que en regresión lineal múltiple): si $x_j$ está correlacionada con una
covariable ausente $x_k$, entonces $\hat\gamma_j$ absorbe parte del efecto de
$x_k$ y el test detecta "no-linealidad" espuria. Esto es coherente con el
supuesto formal ---``no se omiten variables importantes''--- listado en los
tratamientos de diagnóstico logístico (UCLA OARC; University of Guelph).

El test de Wald sobre cada $\hat\gamma_j$, en el modelo aumentado conjunto,
contrasta la linealidad \emph{marginal} de $x_j$ \emph{dado todo lo demás}, que
es la pregunta correcta. Esta versión conjunta (todos los términos $x_j\log x_j$
a la vez) es válida y es la recomendada para modelos multivariados; equivale a
correr un test por variable pero controlando simultáneamente por los demás
términos construidos.

\subsection{Implementación en Python}

\begin{lstlisting}[style=py, caption={Box--Tidwell conjunto con Wald por término y LRT global.}]
import numpy as np
import pandas as pd
import statsmodels.api as sm
import scipy.stats as stats

def box_tidwell_test(df, y, cols_a_testear, cols_control=None):
    """
    H0: gamma_j = 0 (linealidad en el logit) para cada X_j continua > 0.
    - cols_a_testear: continuas positivas -> se agrega X*ln(X).
    - cols_control:   entran como control, NO se testean (negativas, dummies).
    """
    X = pd.DataFrame(index=df.index)
    if cols_control:
        for c in cols_control:
            X[c] = df[c]

    interaction_cols = []
    for c in cols_a_testear:
        v = df[c]
        if not (v > 0).all():
            print(f"[omitida] {c}: tiene valores <= 0")
            continue
        X[c] = v
        X[f"{c}_ln"] = v * np.log(v)
        interaction_cols.append(f"{c}_ln")

    # Eliminar filas con NaN en cualquier columna usada
    X["_y"] = y
    X = X.dropna()
    y_clean = X.pop("_y")

    Xc = sm.add_constant(X)
    m_full = sm.Logit(y_clean, Xc).fit(disp=0)

    print("== Box-Tidwell: Wald por termino (z = gamma/SE ~ N(0,1)) ==")
    resultados = {}
    for col in interaction_cols:
        var = col.replace("_ln", "")
        g, se = m_full.params[col], m_full.bse[col]
        z, p = m_full.tvalues[col], m_full.pvalues[col]
        flag = "*** VIOLACION" if p < 0.05 else "OK"
        print(f"  {var:22s}: gamma={g:+.5f}  z={z:7.3f}  p={p:.4f}  {flag}")
        resultados[var] = {"gamma": g, "se": se, "z": z, "p": p}

    # LRT global: H0 todos los gamma = 0
    m_red = sm.Logit(y_clean, Xc.drop(columns=interaction_cols)).fit(disp=0)
    lr = 2 * (m_full.llf - m_red.llf)
    gl = len(interaction_cols)
    p_lrt = stats.chi2.sf(lr, df=gl)
    print(f"\n== LRT global: G2={lr:.3f}, gl={gl}, p={p_lrt:.4f} ==")

    return m_full, m_red, resultados
\end{lstlisting}

\begin{lstlisting}[style=py, caption={Decisión post-test con criterio de magnitud (clave en $n$ grande).}]
cols_finales = list(cols_control)
for var, r in resultados.items():
    if r["p"] >= 0.05:
        print(f"{var}: OK -> entra sin transformacion adicional")
        cols_finales.append(var)
    elif abs(r["gamma"]) < 1e-2:
        # Significativo pero magnitud trivial: en n grande el Wald
        # detecta violaciones minusculas. Practicamente irrelevante.
        print(f"{var}: p<0.05 pero |gamma|~0 -> violacion irrelevante")
        cols_finales.append(var)
    else:
        # Violacion real -> remedio (termino cuadratico, spline, excluir)
        print(f"{var}: VIOLACION real -> aplicar remedio")
        df[f"{var}_sq"] = df[var] ** 2
        cols_finales += [var, f"{var}_sq"]
\end{lstlisting}

\begin{observacion}[Sensibilidad al tamaño de muestra]
El test de Box--Tidwell \textbf{no es robusto al tamaño muestral}: con $n$
grande, $\operatorname{SE}(\hat\gamma)$ es diminuto y el Wald declara
significativas violaciones de magnitud ínfima. La guía práctica (SPSS,
literatura aplicada) es: con muestras grandes, no preocuparse por una
interacción \emph{apenas} significativa; mirar la \emph{magnitud} de
$\hat\gamma$ junto con el $p$-valor. Con $n\approx 3\times10^5$ esto es esencial.
\end{observacion}

% ==================================================================
\section{Remedios ante violaciones}\label{sec:remedios}
% ==================================================================

Si el test rechaza linealidad con magnitud no trivial, las opciones ordenadas
de menor a mayor complejidad:

\begin{enumerate}[label=\textbf{\alph*.}]
  \item \textbf{Transformación} del predictor (log, raíz, recíproco). Es el
        primer intento natural y suele bastar para variables de escala amplia
        (montos, ingresos).
  \item \textbf{Término polinomial:} añadir $x^2$ (o mayor) al modelo, que
        captura curvatura en el logit. Concuerda con la literatura
        (Hosmer--Lemeshow; notas de UCLA/Guelph).
  \item \textbf{Spline cúbico restringido:} más flexible que un polinomio
        global, con mejor comportamiento en las colas.
  \item \textbf{Excluir la variable} si, además de no-lineal, no aporta poder
        predictivo (confirmar con LRT del modelo con y sin ella).
\end{enumerate}

Tras aplicar un remedio, se \textbf{re-grafica} el logit empírico de la variable
transformada y se \textbf{re-testea} con Box--Tidwell hasta que el supuesto se
sostenga.

\paragraph{Interpretación tras transformar.} La transformación pasa a ser parte
de la \emph{especificación} del modelo: se aplica idénticamente en
entrenamiento, validación, test y producción (idealmente encapsulada en un
\texttt{Pipeline} de \texttt{sklearn}). El coeficiente de $\log(x_j)$ mide el
efecto de multiplicar $x_j$ por $e$ sobre el log-odds; no es comparable en
escala con el de una variable sin transformar.

% ==================================================================
\section{Bondad de ajuste global: test de Hosmer--Lemeshow}\label{sec:hl}
% ==================================================================

\subsection{Por qué no sirve el $\chi^2$ de Pearson clásico}

Para evaluar bondad de ajuste global se compara frecuencias observadas y
esperadas. Con predictores \emph{continuos}, casi cada observación tiene un
patrón de covariables único, generando una tabla de contingencia dispersa
(celdas con conteos $0$ o $1$). En ese régimen, los estadísticos $X^2$ de
Pearson y $G^2$ de devianza \textbf{no} tienen distribución $\chi^2$
aproximada, porque el número de celdas crece con $n$ y las frecuencias
esperadas por celda no crecen. Hosmer y Lemeshow (1980) propusieron agrupar
para restaurar la aproximación.

\subsection{Construcción del estadístico}

\begin{enumerate}[label=\textbf{\arabic*.}]
  \item Ajustar el modelo y calcular las probabilidades ajustadas $\hat\pi_i$
        para toda observación.
  \item Ordenar por $\hat\pi_i$ y partir en $g$ grupos de igual tamaño
        (típicamente $g=10$, \emph{deciles de riesgo}).
  \item En cada grupo $k$: observados $O_k=\sum_{i\in k} y_i$; esperados
        $E_k=\sum_{i\in k}\hat\pi_i = n_k\bar\pi_k$, con $\bar\pi_k$ la
        probabilidad media del grupo.
\end{enumerate}
El estadístico es una suma tipo Pearson sobre éxitos y fracasos de cada grupo:
\begin{equation}
  \hat C = \sum_{k=1}^{g}
  \frac{(O_k - n_k\bar\pi_k)^2}{n_k\,\bar\pi_k(1-\bar\pi_k)}.
  \label{eq:hl}
\end{equation}

\subsection{Distribución y grados de libertad}

\begin{teorema}[Hosmer--Lemeshow, 1980]
Si el modelo logístico es correcto, las frecuencias esperadas por celda son
suficientemente grandes y $g > p+1$, entonces (establecido por simulación)
\[
  \hat C \;\overset{a}{\sim}\; \chi^2_{g-2}.
\]
\end{teorema}

La reducción a $g-2$ grados de libertad (en vez de $g$) se debe a que las
probabilidades ajustadas se estimaron del propio modelo, lo que impone dos
restricciones adicionales (aproximadamente correspondientes al intercepto y a la
pendiente del ajuste sobre los grupos). Un $p$-valor \emph{grande} indica buen
ajuste (no se rechaza que observado $\approx$ esperado). Nótese la lógica
invertida frente a un test de significancia usual: aquí \emph{no} rechazar es lo
deseable.

\begin{observacion}[Limitaciones]
El test de Hosmer--Lemeshow (i) evalúa calibración global, no una forma
específica de mala especificación (p.ej.\ un término cuadrático faltante);
(ii) es sensible al número y método de agrupamiento; (iii) en muestras muy
grandes tiende a rechazar por desviaciones triviales, y en muy pequeñas tiene
poca potencia. Es complementario ---no sustituto--- del diagnóstico de
linealidad. Harrell recomienda además examinar residuos y calibración gráfica.
\end{observacion}

\subsection{¿Cuándo usar Hosmer--Lemeshow? ¿Es opcional?}

Sí, es un test \textbf{opcional y complementario}, no un requisito para que la
regresión logística sea válida. Conviene ubicarlo bien dentro del flujo:

\begin{itemize}
  \item \textbf{Los supuestos que validan el modelo} son los de las Secciones
        \ref{sec:visual}--\ref{sec:remedios}: linealidad en el logit
        (Box--Tidwell), independencia de observaciones, ausencia de
        multicolinealidad (VIF) y tamaño muestral suficiente. Estos son
        \emph{previos} al ajuste y condicionan que la inferencia sobre $\hbeta$
        sea confiable.
  \item \textbf{Hosmer--Lemeshow es \emph{posterior}}: se corre sobre el modelo
        ya ajustado para chequear \emph{calibración global} ---si las
        probabilidades predichas coinciden con las frecuencias observadas por
        grupos de riesgo. No sustituye a los diagnósticos de supuestos.
\end{itemize}

\paragraph{Cuándo sí aporta.}
\begin{itemize}
  \item \textbf{Modelos de predicción de riesgo} donde la probabilidad predicha
        se usa como tal (credit scoring, riesgo clínico, seguros). Ahí la
        calibración es el objetivo mismo: un $\hat\pi=0.1$ debe corresponder a
        $\approx10\%$ de eventos observados. Es su caso de uso canónico y por eso
        aparece en la validación de \emph{scorecards}.
  \item Cuando hay \textbf{predictores continuos} que hacen inaplicable el
        $\chi^2$ de Pearson clásico (Sección \ref{sec:hl}), y se quiere una
        medida de ajuste global resumida en un número.
\end{itemize}

\paragraph{Cuándo es prescindible o engañoso.}
\begin{itemize}
  \item \textbf{Objetivo puramente discriminativo} (ordenar/clasificar, no usar
        la probabilidad como magnitud): interesan más AUC, sensibilidad y
        especificidad. Un modelo puede discriminar bien (AUC alto) y estar mal
        calibrado, o viceversa: son propiedades distintas.
  \item \textbf{Muestras muy grandes} (como el dataset de crédito,
        $n\approx3\times10^5$): el test rechaza por desviaciones triviales
        ---mismo problema de sensibilidad al tamaño que el Box--Tidwell. Con
        $n$ grande conviene apoyarse en \textbf{gráficas de calibración}
        (probabilidad predicha vs observada por bin) en lugar del $p$-valor, que
        casi siempre saldrá significativo.
  \item \textbf{Diagnóstico de una forma concreta de mala especificación:} para
        eso sirven Box--Tidwell (linealidad) o añadir términos y comparar por
        LRT, no Hosmer--Lemeshow, que solo detecta descalibración agregada.
\end{itemize}

En resumen: para el flujo de \emph{este} documento ---verificar linealidad en el
logit--- el test relevante es Box--Tidwell. Hosmer--Lemeshow es un chequeo
\emph{final y opcional} de calibración, especialmente pertinente si el producto
es un modelo de scoring cuyas probabilidades se interpretan literalmente; con
$n$ grande, complementarlo (o reemplazarlo) por una curva de calibración.

\subsection{Implementación en Python}

\begin{lstlisting}[style=py, caption={Test de Hosmer--Lemeshow.}]
import numpy as np
import pandas as pd
import scipy.stats as stats

def hosmer_lemeshow(y_true, y_prob, g=10):
    df = pd.DataFrame({"y": np.asarray(y_true), "p": np.asarray(y_prob)})
    # Deciles de riesgo (igual numero de obs por grupo)
    df["grupo"] = pd.qcut(df["p"], q=g, duplicates="drop")

    agg = df.groupby("grupo", observed=True).agg(
        n=("y", "size"),
        obs=("y", "sum"),
        p_bar=("p", "mean"),
    )
    agg["esp"] = agg["n"] * agg["p_bar"]
    # Estadistico de Hosmer-Lemeshow
    C = (((agg["obs"] - agg["esp"]) ** 2)
         / (agg["n"] * agg["p_bar"] * (1 - agg["p_bar"]))).sum()

    gl = len(agg) - 2
    p_val = stats.chi2.sf(C, df=gl)
    print(f"Hosmer-Lemeshow: C={C:.3f}, gl={gl}, p={p_val:.4f}")
    print("p grande -> no se rechaza buen ajuste (calibracion adecuada)")
    return C, gl, p_val

# Uso tipico tras ajustar el modelo final:
#   m = sm.Logit(y, sm.add_constant(X)).fit()
#   hosmer_lemeshow(y, m.predict(sm.add_constant(X)), g=10)
\end{lstlisting}

% ==================================================================
\section{Flujo de trabajo completo}\label{sec:flujo}
% ==================================================================

Reuniendo todo, el procedimiento para verificar y conseguir la linealidad en el
logit es:

\begin{enumerate}[label=\textbf{\arabic*.}]
  \item \textbf{Clasificar variables:} continuas (candidatas al supuesto),
        binarias/dummies (exentas), y detectar valores anómalos o ceros
        estructurales que requieran limpieza o flags.
  \item \textbf{Visual:} \emph{empirical logit plot} de cada continua. Confirma
        las lineales y señala las sospechosas.
  \item \textbf{Transformar} las visualmente no-lineales (típicamente log en
        variables de escala amplia) y \textbf{re-graficar}.
  \item \textbf{Box--Tidwell conjunto} sobre las continuas positivas del modelo
        final, incluyendo todas las demás como control. Wald por término + LRT
        global.
  \item \textbf{Decidir por variable} combinando $p$-valor y \emph{magnitud} de
        $\hat\gamma$ (crítico en $n$ grande): $p\ge0.05$ o $|\hat\gamma|\approx0$
        $\Rightarrow$ aceptar; violación real $\Rightarrow$ remedio (polinomio,
        spline, o excluir).
  \item \textbf{Iterar} 3--5 hasta que el supuesto se sostenga.
  \item \textbf{Ajustar el modelo final} y verificar bondad de ajuste global con
        Hosmer--Lemeshow, además de multicolinealidad (VIF), independencia y
        tamaño muestral (regla $\geq 10$ eventos por predictor).
  \item \textbf{Inferencia:} Wald/LRT sobre los $\hat\beta$ y odds ratios
        $e^{\hat\beta_j}$ con sus intervalos de confianza.
\end{enumerate}

\begin{center}
\begin{tabular}{lll}
\toprule
Etapa & Estadístico & Distribución bajo $H_0$ \\
\midrule
Coeficientes individuales & Wald $z=\hat\beta_j/\text{SE}$ & $\Nor(0,1)$ \\
Bloque de coeficientes & LRT $G^2$ & $\chi^2_\nu$ \\
Box--Tidwell (por variable) & Wald $z=\hat\gamma/\text{SE}$ & $\Nor(0,1)$ \\
Box--Tidwell (global) & LRT $G^2$ & $\chi^2_{\#\text{términos}}$ \\
Bondad de ajuste & Hosmer--Lemeshow $\hat C$ & $\chi^2_{g-2}$ \\
\bottomrule
\end{tabular}
\end{center}

% ==================================================================
\section{Fuentes y lecturas recomendadas}
% ==================================================================

Recursos con \emph{derivación completa} (no solo enunciado), útiles para
profundizar al nivel de estas notas:

\begin{itemize}
  \item \textbf{G. Rodríguez (Princeton), \emph{Generalized Linear Model
        Notes}, Apéndice A.} Deriva score, información, Newton--Raphson/Fisher
        scoring y los tres tests con ejemplos.
        \url{https://grodri.github.io/glms/notes/a1.pdf}
  \item \textbf{Penn State STAT 504, Lecciones 6--7.} Modelo logístico desde el
        problema de heterocedasticidad de OLS, MLE, Wald y LRT, residuos y
        Hosmer--Lemeshow. \url{https://online.stat.psu.edu/stat504/}
  \item \textbf{MIT OCW 18.650 (Rigollet), Lecturas 21--24, GLMs.} Familia
        exponencial, enlace canónico, información de Fisher, devianza; en video.
  \item \textbf{Shao, \emph{Mathematical Statistics}, Lec.~26} (notas
        U.~Wisconsin): pruebas formales de Wald y Score como $\chi^2$.
        \url{https://pages.stat.wisc.edu/~shao/stat610/stat610-26.pdf}
  \item \textbf{Box \& Tidwell (1962),} \emph{Transformation of the Independent
        Variables}, Technometrics 4:531--550 (fuente primaria de la
        transformación).
  \item \textbf{Pregibon (1980),} \emph{Goodness of Link Tests for GLM}, JRSS-C
        29:15--24 (extensión del test a GLM vía linealización).
  \item \textbf{Hosmer \& Lemeshow (1980)} y Hosmer et al.\ (1997), sobre el
        estadístico de bondad de ajuste y su distribución.
\end{itemize}

\vfill
\begin{center}\small\itshape
Documento generado como material de estudio. Las derivaciones siguen el estándar
de la teoría de verosimilitud; los recursos citados contienen las demostraciones
formales completas.
\end{center}

\end{document}